% @order: 1
% @description: 设置全局字体、页面布局
% @\usepackage[explicit]{titlesec}
% @\usepackage{titletoc}

% !字体和版式设置
% #region

% 页面布局
\geometry{
    left=2.5cm,
    right=2.5cm,
    top=3cm,
    bottom=3cm
}

% 中文字体设置
%\setCJKmainfont{SimSun}[BoldFont=SimHei,ItalicFont=KaiTi]
\setCJKmainfont{Source Han Serif SC}[
    BoldFont={Source Han Sans SC Medium},
    ItalicFont={Kaiti SC}
]
\setCJKsansfont{Source Han Sans SC}[
    BoldFont={Source Han Sans SC Bold}
]
\newCJKfontfamily\heavyfont{Source Han Serif SC Heavy}
\setCJKmonofont{PingFang SC}

% 数学字体设置
\setmathfont{latinmodern-math.otf}
\setmathfont[
    range={"2205},
    Scale=MatchUppercase
]{STIX Two Math}

% 文档分页与数学排版
\let\cleardoublepage\clearpage
\allowdisplaybreaks

% 页面垂直对齐
\raggedbottom

% 章节名称设置
\ctexset{
    chapter/name={},
    chapter/number=\arabic{chapter},
    chapter/format={},
    chapter/aftername={},
    chapter/beforeskip=0pt,
    chapter/afterskip=0pt,
}

% 标题与目录设置
% Chapter样式
\titleformat{\chapter}[display]
    {\raggedright}
    {
        {\fontsize{40pt}{50pt}\selectfont\bfseries C}
        {\fontsize{25pt}{40pt}\selectfont\bfseries \!hapter~}
        {\fontsize{40pt}{50pt}\selectfont\bfseries \thechapter}
    }
    {-15pt}
    {\huge\heavyfont #1}

\titlespacing*{\chapter}
    {0pt}
    {-20pt}
    {30pt}

% Section样式
\tcbset{
    sectionbox/.style={
        enhanced,
        capture=hbox,
        colback=white,
        colframe=black,
        boxrule=0.5pt,
        leftrule=3pt,
        rightrule=3pt,
        arc=5pt,
        left=15pt,
        right=15pt,
        top=6pt,
        bottom=6pt,
        fontupper=\Large\bfseries,
    }
}

\titleformat{\section}
    {\normalfont}
    {}
    {0pt}
    {
        \begin{tcolorbox}[sectionbox]
            \S\thesection\quad #1
        \end{tcolorbox}
    }

\titlespacing{\section}
    {0pt}
    {0pt}
    {-20pt}

% 单位格式
\sisetup{
    per-mode = reciprocal,
    inter-unit-product = \ensuremath{\cdot}
}

% 设置目录
% 目录标题样式
\titleformat{name={\chapter},numberless}[block]
    {\Huge\heavyfont}
    {}
    {0pt}
    {#1}


% 目录条目样式
\titlecontents{chapter}
    [0pt]
    {\addvspace{1em}\bfseries}
    {Chapter~\thecontentslabel\hspace{1em}}
    {}
    {\hfill\contentspage}


% 生成目录
\newcommand{\makecontents}{
    \let\oldcleardoublepage\cleardoublepage
    \let\cleardoublepage\clearpage

    \frontmatter

    \pagestyle{plain}

    \begingroup
        \ctexset{
            contentsname = {目~录}
        }
        \tableofcontents
    \endgroup

    \mainmatter

    \let\cleardoublepage\oldcleardoublepage
}

% 表格列样式设置
\renewcommand{\tabularxcolumn}[1]{m{#1}}
\newcolumntype{L}{>{\hbadness=10000\relax}X}
\newcolumntype{Y}{>{\RaggedRight\arraybackslash}X} % 左对齐可换行
\newcolumntype{Z}{>{\centering\arraybackslash}X}   % 居中对齐

% #endregion
